\section{Thriving in the End Times}

\paragraph{What Is Thinking}


\begin{quotex}
All things that we see standing accomplished in the world are properly the outer material result, the practical realization and embodiment, of Thoughts that dwell in the Great Men sent into the world.
\flright{\textsc{Thomas Carlyle}, \emph{On Heroes, Hero Worship, and the Heroic in History.}}
\end{quotex}


The sophisticated man since the Enlightenment likes to claim that he ``thinks for himself'', without need for any external authority. To say the least, the results of that claim have been quite meagre. Out of the 50,000 thoughts going through your mind each day, how many have been useful, never mind extraordinary or creative?

Ancient Traditional man was less dense, i.e., less sensual, so more detached from the thoughts in his consciousness. He would experience thoughts as commands from the gods.

The ``bright'' ones today go beyond even the Enlightenment ideal and claim that thoughts arise from matter, as the result of bio-electro-chemical processes. This, despite the fact that there is no conceivable explanation of how that might happen. After all, what is the Maxwell equation or chemical formula for any particular thought? If that view is true, then the solution to the problems of the world could be solved by designing the correct pharmaceuticals to produce happy and intelligent thoughts. Keep me in the loop if that happens.

The Great Men of history have been intermediaries, bringing great thoughts into the world. There is no shame in repeating those same thoughts.

\paragraph{Knowledge and Love}


You may have heard it said that all you need is love, and, therefore, knowledge (or gnosis) is unimportant. On the contrary, the two are intimately related. \textbf{Nicholas Cabasilas}, in \emph{The Life in Christ}, explains why:

\begin{quotex}
In fact, it is knowing that causes love and gives birth to it. It is not possible to attain love of anything that is beautiful without first learning how beautiful it is. Since this knowledge is sometimes very ample and complete and at other times imperfect, it follows that the philter of love has a corresponding effect. Some things that are beautiful and good are perfectly known and perfectly loved as befits so great beauty. Others are not clearly evident to those who love them, and love of them is thus more feeble.
\end{quotex}


Hence, gnosis is not an obstacle to love, but rather its prerequisite, and the clearer the knowledge, the stronger the love. He goes on to explain that our knowledge of things is twofold:

\begin{enumerate}


\item That which one may acquire from \textbf{hearsay}, and

\item That which one may learn by \textbf{personal experience}.
\end{enumerate}


In hearsay knowledge, we don't encounter the thing itself, but rather we see it through words about it. This is, at best, to see it inaccurately, since there is nothing in nature that is an adequate copy of it. This uncertain and dim image cannot excite the degree of love proper to it.

When we encounter the things themselves, we gain experience of them. By this experience the very form itself encounters the soul and incites desire, as though it left an imprint corresponding to the good. This is important in one's spiritual life. Cabasilas draws this conclusion:

\begin{quotex}
When therefore love of the Savior produces nothing new or extraordinary in us it proves that we have encountered no more than mere words about Him.
\end{quotex}


In a broader context, there is a lot of talk, ideas, and concepts about God. But that is all hearsay and such knowledge does not lead to love of God. That is because we need to relate concepts of the unknown to what we know by experience, through analogy. But God is unique, so there is nothing else in our experience adequate for that. This cataphatic approach to knowledge of God speaks of an analogy of being, but such analogies are far from precise.

Moreover, thoughts and concepts are always about essences. However, God's essence is the same as his existence, i.e., the act of being or his energies. So we know God's essence through hearsay, but we can know his existence through personal experience. This direct knowing of God's presence is what we mean by gnosis.

\paragraph{The Barriers to Theosis}


\begin{quotex}
Beloved, we are God's children now; what we shall be has not yet been revealed. We do know that when it is revealed we shall be like him, for we shall see him as he is. Everyone who has this hope based on him makes himself pure, as he is pure.
\flright{\textsc{1 John 3}}

Beholding the glory of the Lord we are being changed into the same likeness.
\flright{\textsc{2 Cor 3:18}}
\end{quotex}


Nicholas Cabasilas explains that there are three obstacles to theosis. These are:

\begin{enumerate}


\item \textbf{Nature}: The Divine nature is different from human nature.


\item \textbf{Sin}: A will corrupted by evil separates us from God.


\item \textbf{Death}: In the mortal body, we can see only the dim reflection in the mirror; in this state our bodies are dominated by sense life.
\end{enumerate}


These obstacles are overcome by the following events respectively:

\begin{enumerate}


\item \textbf{Incarnation}: This unites the human and divine natures in one person.

\item \textbf{Crucifixion}: This leads to the forgiveness of sins.

\item \textbf{Resurrection}: This overcomes death and the attraction to sense life.
\end{enumerate}


Cabasilas relates these ideas to the effects of the sacraments, or mysteries, with the aim of salvation. The esoteric path aims beyond this to liberation or \emph{jivanmukti}. That aim is union while still in the mortal body:

\begin{enumerate}


\item Expose our false sense of I, replaced with the mind of Christ.

\item Purify our soul so it becomes the perfect reflector of the Holy Spirit.

\item Move from a life of instinct to a life of intelligence.
\end{enumerate}


\paragraph{Narcissism in Religion}


There is the tendency to confuse knowledge by direct personal experience with profound religious emotions. Of course, such passions can be quite motivational to pursue gnosis, in themselves they fall short of knowledge. Knowledge is a gift of the Holy Spirit, but warm emotional experiences are not among those promised gifts.

What may happen is that the person claims an authority in spiritual matters based solely of one deep experience, particularly if it led to a conversion of some sort. If two such persons confront each other, a narcissistic battle usually ensues about whose experience is deeper or who has the most impressive conversion story.

\paragraph{Strength and Fertility}


\begin{quotex}
The right wing authoritarians want to bring about homogeneity by exclusion and the left wing authoritarians want to bring about homogeneity by inclusion. But they both want to bring about homogeneity. So on the right, that would be more couched in terms of purity, and on the left it would be more couched in terms of equality. 
\flright{\textit{Where Do SJWs Come From?}\footnote{\url{https://www.youtube.com/watch?v=x\_fBYROA7Hk}}}
\end{quotex}


I am always suspicious of such studies that reduce metaphysical and religious issues to the psychological level. The study relates the two wings to rule by fathers and rule by mothers. I'm not surprised that liberals and social justice warriors suffer from anxiety disorders, since their mentality is at odds with the nature of things. Hence its opposite is not the so-called ``right wing authoritarian'', which the study conflates with racism and fascism, but rather the cosmic order.

Strong personal boundaries make for psychological and spiritual health. We see that exemplified in the story of Romulus\footnote{\url{https://gornahoor.net/?page\_id=13248}}, the spiritual father of the races of Western Europe. His first task was to delineate the boundaries of Rome, to determine who is in and who is out, who is excluded and who is included. To do so required the mastery of two beasts, representing strength and fertility: strength to enforce the frontier of the city, and fertility to populate it. Those right wing groups that claim to be saving ``western civilization'' should pay heed to this.

\paragraph{Time Preference}


The notion of Time Preference is crucial to the science of economics and is an area of research in profane psychology. \textbf{Hans-Hermann Hoppe}, in \emph{Democracy: The God that Failed}, explains it this way:

\begin{quotex}
In acting, an actor invariably aims to substitute a more satisfactory for a less satisfactory state of affairs and thus demonstrates a preference for more rather than few goods. Moreover, he invariably considers when in the future his goals will be reached, i.e., the time necessary to accomplish them, as well as a good's duration of serviceability. Thus, he also demonstrates a universal preference for earlier over later goods, and for more over less durable ones. This is the phenomenon of time preference.
\end{quotex}


Children, for example, have a high time preference since they prefer something now rather than later. Studies have shown that children as young as four who have a lower time preference --- e.g., the Stanford Marshmallow Experiment\footnote{\url{https://en.wikipedia.org/wiki/Stanford\_marshmallow\_experiment}} --- are more likely to be successful later in life than children with a higher time preference.

Similarly, societies with a lower time preference will be more successful in the long run since they more carefully prepare for the future. Hoppe makes the case that monarchies have a lower time preference than democracies; hence the latter are doomed since they will tend to squander their capital.

Historically, civilizations with a strong spiritual base have thrived for much longer durations. The religious mind, since it takes into account postmortem states, will have the lowest time preferences of all. They also feel a spiritual kinship with their ancestors and descendants, a feeling that secular societies lack.

The ancient Romans and Greeks, for example, kept a hearth fire burning for their ancestors and they expected future generations to do the same for them. They are a dim image, as though through a glass darkly, of postmortem states. The Middle Ages developed that notion: the Church Suffering required the prayers of the living and the Church Triumphant would then intercede for the living. The effects of the increasing lack of awareness for this multi-generational connection are obvious.


\flrightit{Posted on 2016-11-03 by Cologero}

\begin{center}* * *\end{center}

\begin{footnotesize}
\begin{sffamily}
\texttt{Ann K. on 2016-11-04 at 08:23 said:}

Beautiful! Thank you.

\hfill

\texttt{Tom Rowsell on 2016-11-04 at 15:53 said:}

Interesting final thought -- that prayers for the living and the dead are a development of our old pagan customs of ancestor worship. Each facilitates the continuity of tradition.

\hfill

\texttt{G. Eiríksson on 2016-11-16 at 09:44 said:}

Best to think not at all.


\hfill

\end{sffamily}
\end{footnotesize}

